% ============================== CI-3 ==============================
\reading{CI-3}{The Global Drivers of Private Credit}{STRIKE FIRST}{2}

\begin{cikeybox}[title={The three objectives, in the spine's own words}]
\begin{enumerate}[label=\textbf{\alph*.}]
\item Describe the structure and characteristics of private credit funds, as well as the
patterns in the global growth of the asset class.
\item Discuss the potential drivers of private credit growth and interpret regression results
about their impact.
\item Compare private credit loans with bank loans and assess how changes in cost of equity,
cost of debt, and leverage have influenced the competitive dynamics between them.
\end{enumerate}
\greyline{Source: ``The Global Drivers of Private Credit'', BIS (February 2025). The class
notes label this Reading 102; it is CI-3 on the 2026 spine. Objective (b) contains the one
substantive correction in this chapter, logged at the back.}
\end{cikeybox}

% ------------------------------------------------------------------
\lo{CI-3 a}{Describe the structure and characteristics of private credit funds, as well as the patterns in the global growth of the asset class.}

\subsection*{Structure and characteristics of private credit funds (PCFs)}

Private credit funds (PCFs) provide loans directly to borrowers such as small- and
medium-sized nonfinancial firms, and these loans are generally held for their full term. PCFs
are usually structured as \term{closed-end funds}, so investor capital remains committed for
around five to eight years. This reduces liquidity and maturity transformation risk, but
because these funds are unlisted, they are relatively illiquid.

In the United States, \term{business development companies (BDCs)} make up roughly one-fifth
of the private credit market. Unlike traditional PCFs, BDCs are often publicly listed and
designed to appeal more to retail investors by allowing relatively earlier redemptions
through open-end features.

The investor base of PCFs is still mainly institutional, since these investors can lock in
capital for longer periods. However, retail participation is gradually rising. This matters
because if PCFs increasingly depend on retail money, they may begin to face asset-liability
mismatch risks similar to those faced by banks, which rely heavily on short-term funding.

\begin{cifigblock}
{\centering
\begin{tikzpicture}[
  font=\sffamily\scriptsize,
  inv/.style={draw=FRMRule, fill=PaleGreen, text width=1.95cm, align=center,
              minimum height=1.05cm, inner sep=3pt, rounded corners=1pt},
  int/.style={draw=FRMRule, fill=PaleRed, text width=1.95cm, align=center,
              minimum height=1.15cm, inner sep=3pt, rounded corners=1pt},
  wide/.style={draw=FRMRule, fill=PaleNavy, text width=8.9cm, align=center,
               minimum height=0.85cm, inner sep=3pt, rounded corners=1pt},
  bank/.style={draw=FRMRule, fill=FRMSoft, text width=1.5cm, align=center,
               minimum height=1.05cm, inner sep=3pt, rounded corners=1pt},
  lbl/.style={font=\sffamily\bfseries\scriptsize\color{FRMNavy}, align=left},
  fl/.style={-{Stealth[length=4pt]}, FRMNavy, line width=0.7pt}
]
\node[inv] (pf) at (0,0) {\textbf{PFs}\\ pension funds};
\node[inv, right=3mm of pf] (ic) {\textbf{ICs}\\ insurance corporations};
\node[inv, right=3mm of ic] (swf) {\textbf{SWFs}\\ sovereign wealth funds};
\node[inv, right=3mm of swf] (ret) {\textbf{Retail}};
\node[lbl, left=4mm of pf, text width=2.1cm] {End\\investors};

\node[int] (f1) at (0,-2.15) {\textbf{Funds}\\ closed-ended, GP/LP structure};
\node[int, right=3mm of f1] (f2) {\textbf{Funds}\\ open-ended};
\node[int, right=3mm of f2] (f3) {\textbf{BDCs}\\ business development companies};
\node[int, right=3mm of f3] (f4) {\textbf{CLOs}\\ collateralised loan obligations};
\node[lbl, left=4mm of f1, text width=2.1cm] {Inter\-mediaries};

\node[bank, right=9mm of f4] (bk) {\textbf{Banks}};
\draw[fl, FRMGold] (bk.north) |- ($(ret.east)+(0.4,0)$) -- (ret.east)
  node[midway, above, font=\sffamily\scriptsize\color{FRMGold}, fill=white, inner sep=1pt] {};
\node[font=\sffamily\scriptsize\color{FRMGold}, fill=white, inner sep=1.5pt]
  at ($(bk.north)+(0,0.55)$) {leverage};

\node[wide] (terms) at (3.0,-3.6) {floating rate loans \ \textbullet\ \ covenant heavy /
  heavier \ \textbullet\ \ degree of control \ \textbullet\ \ customised terms \ \textbullet\
  \ few lenders};
\node[wide, fill=PaleGold] (corp) at (3.0,-4.95) {\textbf{Corporates}: small and
  medium-sized firms, some highly leveraged};
\node[lbl, text width=2.1cm, anchor=east] at ($(corp.west)+(-0.4,0)$)
  {Borrowers};

\draw[fl] (pf.south)  -- (f1.north);
\draw[fl] (ic.south)  -- (f2.north);
\draw[fl] (swf.south) -- (f3.north);
\draw[fl] (ret.south) -- (f4.north);
\draw[fl] (f1.south) -- ([xshift=-2cm]terms.north);
\draw[fl] (f2.south) -- ([xshift=-0.7cm]terms.north);
\draw[fl] (f3.south) -- ([xshift=0.7cm]terms.north);
\draw[fl] (f4.south) -- ([xshift=2cm]terms.north);
\draw[fl] (terms.south) -- (corp.north);
\end{tikzpicture}\par}
\figcap{Figure CI-3.1. The intermediation chain, redrawn from the notes' IMF figure. PCFs sit
between end investors and middle-market corporate borrowers, and the middle band is what
makes the loans different from bank loans: floating rate, covenant heavy, few lenders, high
degree of control. Banks appear only at the edge, supplying leverage to the intermediaries
rather than lending to the borrowers, which is the structural reason banks' exposure to this
market is indirect. Source: IMF staff. GP = general partners, LP = limited partners.}
\end{cifigblock}

\subsection*{Lending strategies}

PCFs typically follow a \term{focused and specialised lending approach}, targeting borrowers
that may not have access to traditional bank financing. Common strategies include:

\begin{center}
\renewcommand{\arraystretch}{1.25}
\begin{tabularx}{\textwidth}{@{}p{3.4cm} >{\raggedright\arraybackslash}X@{}}
\toprule
\rowcolor{FRMSoft}
{\sffamily\bfseries\color{FRMNavy}Strategy} &
{\sffamily\bfseries\color{FRMNavy}What it is}\\
\midrule
\term{Direct lending} & Core strategy, often with floating interest rates and strong
covenants\\
\term{Mezzanine lending} & Subordinated debt with possible equity upside\\
\term{Asset-based lending} & Loans backed by collateral\\
\term{Distressed debt investing} & Buying high-risk debt at deep discounts\\
\bottomrule
\end{tabularx}
\end{center}

\subsection*{Global growth of private credit}

\begin{itemize}
\item Private credit has expanded rapidly, with AUM rising from \$0.2T to \$2.5T over two
decades.
\item The U.S. dominates (90\%), while Europe, the UK, and Canada make up the rest.
\item The sector has also broadened beyond traditional industries like manufacturing and
technology into areas such as cleantech and life sciences.
\end{itemize}

\subsection*{Portfolio concentration, risks and tradeoff}

PCFs are highly concentrated (measured by HHI) and less diversified than banks. This is
intentional, as fund managers build sector expertise to lend to riskier borrowers (high
leverage, low collateral) and offer customised loans.

This model creates a clear tradeoff. While specialisation can generate higher returns, it
also exposes funds to sector-specific downturns. Institutional investors are better able to
manage this risk through diversification, whereas retail investors may underestimate it.

% ------------------------------------------------------------------
\lo{CI-3 b}{Discuss the potential drivers of private credit growth and interpret regression results about their impact.}

\subsection*{Potential growth drivers}

The growth of private credit funds (PCFs) is mainly driven by four structural factors:

\begin{enumerate}
\item \term{Monetary policy.} Low policy rates compress banks' margins (assets fall faster
than liabilities), reducing their willingness to lend.
  \begin{itemize}
  \item PCFs benefit because their assets and liabilities are mostly floating-rate, allowing
  them to earn relatively higher returns and attract investors.
  \end{itemize}
\item \term{Banking sector efficiency (FI index).} Banks often avoid lending to SMEs due to
weak collateral and transparency.
  \begin{itemize}
  \item PCFs step in by offering flexible lending with customised covenants instead of strict
  collateral requirements.
  \end{itemize}
\item \term{Regulatory stringency.} Higher capital requirements and stricter oversight make
banks more cautious.
  \begin{itemize}
  \item Riskier borrowers shift toward PCFs, which operate with more flexibility.
  \end{itemize}
\item \term{Corporate leverage.} As non-financial firms take on more debt, PCFs have shown a
higher appetite for highly leveraged borrowers compared to traditional lenders.
  \begin{itemize}
  \item PCFs are more willing than banks to lend to highly leveraged firms.
  \end{itemize}
\end{enumerate}

\subsection*{Regression analysis}

The BIS analysis quantified the impact of these drivers across 45 countries. The results
demonstrate that banking inefficiency and interest rates are the most significant predictors
of private credit expansion.

\begin{center}
\renewcommand{\arraystretch}{1.25}
\small
\begin{tabularx}{\textwidth}{@{}p{3.5cm} p{3.0cm} p{2.5cm}
                              >{\raggedright\arraybackslash}X p{1.7cm}@{}}
\toprule
\rowcolor{FRMSoft}
{\sffamily\bfseries\color{FRMNavy}Variable} &
{\sffamily\bfseries\color{FRMNavy}Coefficient type} &
{\sffamily\bfseries\color{FRMNavy}Direction of the 1 SD change} &
{\sffamily\bfseries\color{FRMNavy}Effect on private credit activity} &
{\sffamily\bfseries\color{FRMNavy}Influence level}\\
\midrule
FI Index (bank efficiency) & Negative ($-1.68$ to $-1.91$) &
\textcolor{FRMRed}{\bfseries decline} & 33\% increase & Highest\\
Policy rate & Negative ($-0.03$) &
\textcolor{FRMRed}{\bfseries decline} & 12\% increase & High\\
Regulatory stringency & Positive (0.46 to 0.67) &
\textcolor{FRMGreen}{\bfseries increase} & 2\%--4\% increase & Moderate\\
Corporate leverage & Positive (0.001) &
\textcolor{FRMGreen}{\bfseries increase} & under 1\% increase & Minimal\\
\bottomrule
\end{tabularx}
\end{center}
\figcap{Table CI-3.1. The third column is new and is the correction. The notes' version of
this table has four columns and reads ``Impact of 1 Std.\ Deviation Change: 33\% Increase in
PC activity'' against a negative coefficient, which as printed says that a \emph{rise} in
bank efficiency raises private credit. It does not. A negative coefficient means the increase
follows a \emph{fall} in the variable.}

\begin{citrapbox}[title={The correction, stated plainly}]
Two of the four rows are direction-reversed as printed in the notes. A one-standard-deviation
\textbf{decline} in the FI index is associated with roughly a 33\% increase in private credit
activity, and a one-standard-deviation \textbf{decline} in the policy rate with roughly a 12\%
increase. The other two rows are increases and are correct as printed, which is exactly what
makes the error easy to miss: the table looks internally consistent.\par\vspace{4pt}
The sign of the coefficient and the direction of the move always have to be read together.
Read the whole row: negative coefficient, so the effect column describes what happens when
the variable \emph{falls}.\par\vspace{4pt}
{\footnotesize\color{FRMGrey}Resolved against SchweserNotes Book 5, LO 102.b. \schweserbadge}
\end{citrapbox}

\begin{cifigblock}
{\centering
\begin{tikzpicture}
\begin{axis}[
  width=12.0cm, height=5.0cm,
  xbar, bar width=13pt,
  xmin=0, xmax=42,
  xlabel={\sffamily\scriptsize Increase in private credit activity per 1 SD move in the stated direction (\%)},
  symbolic y coords={Corporate leverage, Regulatory stringency, Policy rate, FI Index},
  ytick={Corporate leverage, Regulatory stringency, Policy rate, FI Index},
  bar shift=0pt,
  tick label style={font=\sffamily\scriptsize},
  label style={font=\sffamily\scriptsize},
  nodes near coords,
  nodes near coords style={font=\sffamily\scriptsize\color{FRMNavy},
     fill=white, inner sep=1.5pt, anchor=west, xshift=3pt},
  point meta=explicit symbolic,
  xmajorgrids, grid style={FRMRule, dashed},
  axis line style={FRMRule},
  enlarge y limits=0.22
]
\addplot[fill=FRMRed!65, draw=FRMRed, bar shift=0pt] coordinates {
  (33,{FI Index})   [33\% when the index FALLS]
  (12,{Policy rate})[12\% when the rate FALLS]
};
\addplot[fill=FRMGreen!65, draw=FRMGreen, bar shift=0pt] coordinates {
  (4,{Regulatory stringency}) [2--4\% when stringency RISES]
  (0.8,{Corporate leverage})  [under 1\% when leverage RISES]
};
\end{axis}
\end{tikzpicture}\par}
\figcap{Figure CI-3.2. The same four results with the direction carried into the colour: red
bars are effects that follow a \emph{fall} in the variable, green bars follow a \emph{rise}.
The ranking is the examinable part, and it is very lopsided. The FI index alone is worth
almost three times the policy rate and roughly ten times regulatory stringency, while
corporate leverage is statistically present and economically almost irrelevant. The
regulatory stringency bar is drawn at the top of its stated 2\% to 4\% range.}
\end{cifigblock}

% ------------------------------------------------------------------
\lo{CI-3 c}{Compare private credit loans with bank loans and assess how changes in cost of equity, cost of debt, and leverage have influenced the competitive dynamics between them.}

Private credit loans are more flexible and cater to riskier borrowers who may not qualify for
bank loans, whereas banks follow stricter lending standards and prefer safer borrowers.

\begin{cidefbox}[title={Cost of capital framework}]
This mainly depends on \term{cost of equity}, \term{cost of debt}, and \term{leverage} (use of
debt/equity).
\end{cidefbox}

After the 2008 crisis, banks faced a higher cost of equity due to regulation and risk, forcing
them to rely more on debt. In contrast, BDCs had a lower cost of equity and started increasing
their leverage by using more debt instead of equity. Even though BDCs have a slightly higher
cost of debt than banks, debt is still cheaper than equity, so this shift helped reduce their
overall cost of capital.

\begin{cifigblock}
{\centering
\begin{tikzpicture}[
  font=\sffamily\scriptsize,
  hd/.style={draw=FRMRule, fill=FRMSoft, text width=6.5cm, align=center,
             minimum height=0.7cm, inner sep=3pt, rounded corners=1pt},
  cell/.style={draw=FRMRule, fill=white, text width=3.0cm, align=center,
               minimum height=0.95cm, inner sep=3pt, rounded corners=1pt},
  note/.style={text width=6.5cm, align=center, font=\sffamily\scriptsize\itshape\color{FRMGrey}},
  res/.style={draw=FRMRule, fill=PaleGold, text width=6.5cm, align=center,
              minimum height=0.72cm, inner sep=3pt, rounded corners=1pt},
  dn/.style={-{Stealth[length=4pt]}, FRMNavy, line width=0.7pt}
]
% left: bank
\node[hd] (bh) at (0,0) {\bfseries\color{FRMNavy}WACC of a BANK};
\node[cell] (b1) at (-1.75,-1.05) {cost of equity\\ \textcolor{FRMRed}{\bfseries higher}};
\node[cell] (b2) at ( 1.75,-1.05) {cost of debt\\ \textcolor{FRMGreen}{\bfseries low}};
\node[note] (bn) at (0,-2.15) {higher CoE from regulation and risk perception;\\
  cheap deposits keep CoD low};
\node[res] (br) at (0,-3.05) {\bfseries funding advantage, historically};
\draw[dn] (bh.south) -- ++(0,-0.12) -| (b1.north);
\draw[dn] (bh.south) -- ++(0,-0.12) -| (b2.north);
\draw[dn] (bn.south) -- (br.north);

% right: BDC
\node[hd] (dh) at (7.6,0) {\bfseries\color{FRMNavy}WACC of a BDC};
\node[cell] (d1) at (5.85,-1.05) {cost of equity\\ \textcolor{FRMGreen}{\bfseries lower}};
\node[cell] (d2) at (9.35,-1.05) {cost of debt\\ \textcolor{FRMRed}{\bfseries slightly higher}};
\node[note] (dn2) at (7.6,-2.15) {lower CoE post-2008;\\ uses more debt instead of equity};
\node[res] (dr) at (7.6,-3.05) {\bfseries reduced WACC through leverage};
\draw[dn] (dh.south) -- ++(0,-0.12) -| (d1.north);
\draw[dn] (dh.south) -- ++(0,-0.12) -| (d2.north);
\draw[dn] (dn2.south) -- (dr.north);
\end{tikzpicture}\par}
\figcap{Figure CI-3.3. Side by side, the point is that the BDC wins on the \emph{expensive}
leg and loses on the cheap one. Its cost of debt is the higher of the two, but debt is
cheaper than equity for both parties, so the lower cost of equity plus more leverage beats the
bank's deposit funding advantage. The notes say this directly: cost of debt played only a
minor role in the shift.}
\end{cifigblock}

Over time, this reduced the funding advantage that banks historically had. As BDCs lowered
their cost of capital, they could offer more competitive loan pricing, leading to the growth
of private credit.

The main drivers of this shift were \term{lower cost of equity} and \term{higher leverage} for
BDCs, while cost of debt played only a minor role.

\begin{citrapbox}[title={The condition under which the advantage reverses}]
In a rising interest rate environment, BDCs may lose some of this advantage because their cost
of debt increases faster than that of banks. The whole competitive story in this objective is
therefore rate-dependent, and the direction of rates is the switch.
\end{citrapbox}
